% Context from chapters-tex/shannon.tex; for mathematical reference, not compilation.
tension to the complex plane.
 
Shannon's 1949 paper placed the sampling theorem at the heart of communication theory, while acknowledging earlier interpolation results and Nyquist's work on transmission rates.
 
Figure~\ref{fig-sampling-aliasing} illustrates the reconstruction argument and the aliasing that occurs when spectral replicas overlap; see also Figure~\ref{fig-aliasing}.

\begin{figure}[tbp]
\centering
\BookDrawing{tikz/shannon-sampling-aliasing}
\caption{Sampling and Shannon interpolation in the frequency and time domains, with unit sample spacing. Left: the bandlimiting condition of Theorem~\ref{thm-shannon-sampling} holds. Right: overlapping spectral replicas cause aliasing; the reconstructed signal still passes through every sample. The rows show the original signal, sampling, the interpolation kernel, and reconstruction.}
\label{fig-sampling-aliasing}
\end{figure}





\begin{thm} \label{thm-shannon-sampling}
	If $f\in C(\RR)$, $|f(x)| \leq C(1+|x|)^{-3}$ for some $C$ and $\supp(\hat f) \subset [-\pi/s,\pi/s]$, then one has
	\eql{\label{eq-shannong-interp}
		\foralls x \in \RR, \quad 
		f(x) = \sum_n f(n s) \sinc(x/s-n) \qwhereq
		\sinc(u) = \frac{\sin(\pi u)}{\pi u},\quad \sinc(0)=1
	}
	with uniform convergence.
\end{thm}

\begin{proof} 
	Set $g \eqdef f(s \cdot)$. Then $\hat g=1/s \hat f(\cdot/s)$, as follows from the substitution $z=s x$:
	\eq{
		\hat g(\om) = \int f(s x) e^{-\imath \om x} \d x = \frac{1}{s} \int f(z) e^{-\imath (\om/s) z} \d z = \hat f(\om/s)/s, 
	} 
	It is therefore enough to prove the result for unit sample spacing:
	$$
		g(x) = \sum_n g(n) \sinc(x-n), 
	$$
	For the rest of the proof, write $f$ in place of $g$.
	 
	The compact support hypothesis implies $\hat f(\om) = 1_{[-\pi,\pi]}(\om) \hat f_P(\om)$.  
	Combining the inversion formula~\eqref{eq-i-ft} with the Poisson formula~\eqref{eq-poisson-formula} gives
	\eq{
		f(x) = \frac{1}{2\pi} \int_{-\pi}^\pi \hat f_P(\om) e^{\imath \om x} \d \om
		= \frac{1}{2\pi} \int_{-\pi}^\pi \sum_n f(n) e^{\imath \om (x-n)} \d \om.
	} 
	The decay of $f$ gives $\int_{-\pi}^\pi \sum_n |f(n) e^{\imath \om (x-n)}| \d \om = 2\pi\sum_n |f(n)| < +\infty$. We may therefore exchange summation and integration:
	\eq{
		f(x) = \sum_n f(n)  \frac{1}{2\pi} \int_{-\pi}^\pi e^{\imath \om (x-n)} \d \om = \sum_n f(n) \sinc(x-n).
	}
The series converges uniformly by the Weierstrass test, since $|\sinc(x-n)|\leq 1$ and $\sum_n|f(n)|<\infty$.
\end{proof}

\begin{figure}[tbp]
\centering
\BookDrawing{tikz/shannon-sinc}
\caption{The sinc interpolation kernel.}
\label{fig-review-shannon-sinc}
\end{figure}
The sinc kernel decays slowly and oscillates, making it inconvenient for local interpolation. Oversampling gives more flexibility: if $\supp(\hat f)\subset[-\pi/s',\pi/s']$ with $s'>s$, the spectrum occupies a strict subinterval of the Nyquist interval. One can then choose a reconstruction kernel with a smooth, compactly supported Fourier transform. The resulting kernel decays faster than every inverse power of time. Smooth compact spectral support does not, in general, imply exponential decay.

 
For local interpolation, cardinal B-splines offer another useful family of kernels. Define $\phi_0=1_{[-1/2,1/2]}$ and $\phi_k = \phi_{k-1} \star \phi_0$. The function $\phi_k$ is piecewise polynomial of degree $k$, with support $[-(k+1)/2,(k+1)/2]$. For $k\geq1$, it belongs to $C^{k-1}$ and has a bounded weak derivative of order $k$.
The reconstruction takes the form $f \approx \tilde f \eqdef \sum_n a_n \phi(\cdot-n)$, with $\phi=\phi_k$. The interpolation conditions $\tilde f(n)=f(n)$ define a linear system for the coefficients $(a_n)_n$ in terms of the samples $(f(n))_n$. For general sample data, the identity $a_n=f(n)$ holds only for $k \in \{0,1\}$, corresponding to piecewise constant and piecewise affine interpolation.
 
A common choice is cubic spline interpolation, which corresponds to $k=3$.

\texttt{Associated code: test\_sampling.m}

\begin{figure}[tbp]
\centering
\BookDrawing{tikz/shannon-splines}
\caption{Cardinal splines as basis functions for interpolation.}
\label{fig-sampling-splines}
\end{figure}



When the spectral replicas overlap, sampling creates aliasing: high-frequency components become indistinguishable from lower-frequency components. A low-pass filter applied before sampling attenuates frequencies beyond the Nyquist interval $[-\pi/s,\pi/s]$. This reduces aliasing at the cost of losing high-frequency information; see Figure~\ref{fig-aliasing}.

\begin{figure}[tbp]
\centering
\BookDrawing{tikz/shannon-aliasing}
\caption{Aliasing of a sine wave. This periodic example is interpreted through its discrete spectral lines rather than as an $L^2(\RR)$ signal.}
\label{fig-aliasing}
\end{figure}

   
\paragraph{Quantization.}

To store or transmit a sampled signal digitally, one quantizes its values to finite precision.
  
Section~\ref{sec-transform-coding} develops transform coding, in which an invertible linear transform, usually orthogonal, is applied before quantization. Transforming the data can make the coefficients easier to compress. Directly quantizing the sampled values is a simpler alternative.

For example, a step size $s=1/N$ gives the finite vector $(u_n \eqdef f(n/N))_{n=1}^{N} \in \RR^N$. This records only finitely many samples, which do not determine a general bandlimited signal. Exact recovery from these data requires additional assumptions, such as a specified finite-dimensional signal model.

\begin{figure}[tbp]
\centering
\BookDrawing{tikz/shannon-quantizer}
\caption{A uniform scalar quantizer.}
\label{fig-review-shannon-quantizer}
\end{figure}

For a quantization step $T>0$, the quantizer $v_n = Q_T(u_n) \in \ZZ$ records the index of the nearest multiple of $T$, i.e.
\eq{
	v = Q_T(u) \quad\Leftrightarrow\quad
	v-\frac{1}{2} \leq u/T < v+\frac{1}{2},
}
as illustrated by the quantizer diagram. Dequantization replaces each index by the cell midpoint $D_T(v) \eqdef T v$, which minimizes the worst-case error within that cell. The resulting error is bounded by $T/2$: $|D_T(Q_T(u))-u| \leq T/2$.
 
Quantization introduces the loss relative to the sampled vector. Source coding and decoding then preserve the quantized symbols exactly.


\myfigure{
\includegraphics[width=.8\linewidth]{1-shannon/quantize/quantization}
            
                                                                     
                                                                     
                                                                     
                                                                       
                 
                 
                 
                 
  
}{Quantizing an image with $K\in\{2,3,4,16\}$ gray levels. The nor